\section{Conceptual Foundations: Tasks, Assistance, Modality, and Analytic Agency}
\label{sec:foundations}
\draftnote{Morris. Condense submitted background and task rationale. Establish terms here so later sections compare evidence without repeating definitions.}
\subsection{Relational and Multiscale Life-Science Analysis}
\plan{Distinguish explicit networks, derived relationships, spatial neighborhoods, and multiscale structures. Explain how representation and aggregation mediate biological interpretation. Retain concise examples from the submitted background rather than surveying all bioimage analysis.}
\subsection{Analytic Tasks as Human Intent}
\plan{Ground the five categories in prior task frameworks and explain overlap and abstraction level. Show a workflow moving among tasks. Present a task synthesis with a mapping to prior frameworks, rather than a claim that existing taxonomies are obsolete.}
\subsection{Computational Assistance and Its Overlapping Modes}
\plan{Define assistance by its role within interactive analysis. Distinguish Algorithmic transformation, Adaptive context-dependent intervention, Conversational interpretation through language, and Immersive computational use of spatial or embodied context. Include boundary cases and multilabel examples.}
\subsection{Visualization Modality as Interaction Environment}
\plan{Define Desktop/Planar, Large Display, VR, AR/MR, and CAVE consistently with the protocol. Explain how hybrid environments are recorded. Separate display properties, input channels, and computational assistance.}
\subsection{Locus of Analytic Agency}
\plan{Develop who initiates, determines scope, executes transformations, interprets outputs, validates claims, and can intervene or reverse decisions. Use one comparison workflow to show changing responsibility across stages. Distinguish automated execution, initiative, and scientific authority. Ground this in verified mixed-initiative, automation, and shared-task sources from the agreed reading set.}
\subsection{Positioning Against Existing Reviews and Frameworks}
\plan{Compare biological/network surveys, immersive reviews, and human--AI frameworks by domain, task coverage, assistance, modality, and evaluation. Identify the relationship this STAR adds and use that evidence to support the Introduction gap. Do not claim prior work ignores agency. Planned artifacts: task-framework mapping, survey-comparison table, and a worked agency example.}
\draftnote{Recover and verify the exact Horvitz, Domova/Vrotsou, Holter/El-Assady, and Monadjemi foundational works before adding citations. D/M/Target remains supplemental.}
